\begin{table}[htbp]
	\centering
	\caption{Auszug gängiger KI-Video-Generatoren zum 25.01.2026 (Eigene Darstellung)}
	\label{tab:ki_video_generatoren}
	\begin{threeparttable}
		% 4 Spalten: Modell, Anbieter und Plattform sind flexibel (X), 
		% die letzte Spalte ist zentriert (c)
		\begin{tabularx}{\textwidth}{>{\raggedright\arraybackslash}X >{\raggedright\arraybackslash}X >{\raggedright\arraybackslash}X c}
			\toprule
			\textbf{Modell (Version)} & \textbf{Anbieter} & \textbf{Plattform(en)} & \textbf{Text-/Image-to-Video} \\
			\midrule
			Veo (3.1) & Google DeepMind & Gemini, Perplexity & ja / ja \\
			Sora (2.0) & OpenAI & ChatGPT, Bing & ja / ja \\
			Pika (2.5) & Pika Labs & Pika Art & ja / ja \\
			Runway (4.5) & Runway AI & Runway & ja / ja \\
			Kling AI Video (2.6) & Kuaishou & Klingai & ja / ja \\
			Luma Dream Machine Ray (3.0) & Luma AI & Luma AI & ja / ja \\
			Imagine (0.9) & xAI & Grok & ja / ja \\
			\bottomrule
		\end{tabularx}
	\end{threeparttable}
\end{table}